\documentclass[11pt]{article}

\usepackage[a4paper,margin=1in]{geometry}
\usepackage[T1]{fontenc}
\usepackage[utf8]{inputenc}
\usepackage{libertinus}
\usepackage{microtype}
\usepackage{amsmath,amssymb,amsthm,mathtools}
\usepackage{mathrsfs}
\usepackage{booktabs,longtable,tabularx,array}
\usepackage{enumitem}
\usepackage{xcolor}
\usepackage{graphicx}
\usepackage{listings}
\usepackage[most]{tcolorbox}
\usepackage{hyperref}
\usepackage[nameinlink,noabbrev]{cleveref}
\usepackage{fancyhdr}
\usepackage{lastpage}

\definecolor{deepblue}{HTML}{17365D}
\definecolor{teal}{HTML}{0B6E69}
\definecolor{darkgreen}{HTML}{226622}
\definecolor{amber}{HTML}{A05A00}
\definecolor{brick}{HTML}{8C2D24}
\definecolor{softblue}{HTML}{EEF4FA}
\definecolor{softgreen}{HTML}{EEF8F1}
\definecolor{softamber}{HTML}{FFF5E6}
\definecolor{softred}{HTML}{FBEDEC}
\definecolor{softgray}{HTML}{F4F5F7}

\hypersetup{
  colorlinks=true,
  linkcolor=deepblue,
  citecolor=teal,
  urlcolor=teal,
  pdftitle={Kummer--Chebotarev Progress on the Alaoglu--Erdos Conjecture},
  pdfauthor={OpenAI GPT-5.6 Pro},
  pdfsubject={Alaoglu--Erdos conjecture; Kummer theory; Chebotarev density; algebraic K-theory}
}

\pagestyle{fancy}
\fancyhf{}
\fancyhead[L]{Kummer--Chebotarev progress on Alaoglu--Erd\H{o}s}
\fancyhead[R]{22 August 2026}
\fancyfoot[C]{\thepage\ of \pageref{LastPage}}
\setlength{\headheight}{14pt}

\setlist[itemize]{leftmargin=1.5em,itemsep=0.25em,topsep=0.35em}
\setlist[enumerate]{leftmargin=1.7em,itemsep=0.25em,topsep=0.35em}

\numberwithin{equation}{section}
\newtheorem{theorem}{Theorem}[section]
\newtheorem{proposition}[theorem]{Proposition}
\newtheorem{lemma}[theorem]{Lemma}
\newtheorem{corollary}[theorem]{Corollary}
\newtheorem{criterion}[theorem]{Criterion}
\theoremstyle{definition}
\newtheorem{definition}[theorem]{Definition}
\newtheorem{experiment}[theorem]{Computation}
\theoremstyle{remark}
\newtheorem{remark}[theorem]{Remark}
\newtheorem{warning}[theorem]{Warning}

\newcommand{\Q}{\mathbb Q}
\newcommand{\Z}{\mathbb Z}
\newcommand{\N}{\mathbb N}
\newcommand{\F}{\mathbb F}
\newcommand{\Gm}{\mathbb G_m}
\newcommand{\Gal}{\operatorname{Gal}}
\newcommand{\Span}{\operatorname{span}}
\newcommand{\rank}{\operatorname{rank}}
\newcommand{\cont}{\operatorname{cont}}
\newcommand{\ord}{\operatorname{ord}}
\newcommand{\disc}{\operatorname{disc}}
\newcommand{\Sym}{\operatorname{Sym}}
\newcommand{\Frob}{\operatorname{Frob}}
\newcommand{\dens}{\operatorname{dens}}
\newcommand{\Res}{\operatorname{Res}}
\newcommand{\Supp}{\operatorname{Supp}}
\newcommand{\vpp}[2]{v_{#1}(#2)}
\newcommand{\ReportA}{\textit{Unified Report A}}
\newcommand{\abs}[1]{\lvert#1\rvert}
\newcommand{\norm}[1]{\lVert#1\rVert}

\newtcolorbox{resultbox}[1][]{
  enhanced,breakable,colback=softgreen,colframe=darkgreen,
  boxrule=0.7pt,arc=1.5mm,left=2mm,right=2mm,top=1.5mm,bottom=1.5mm,
  title={#1},fonttitle=\bfseries
}
\newtcolorbox{statusbox}[1][]{
  enhanced,breakable,colback=softblue,colframe=deepblue,
  boxrule=0.7pt,arc=1.5mm,left=2mm,right=2mm,top=1.5mm,bottom=1.5mm,
  title={#1},fonttitle=\bfseries
}
\newtcolorbox{gapbox}[1][]{
  enhanced,breakable,colback=softamber,colframe=amber,
  boxrule=0.7pt,arc=1.5mm,left=2mm,right=2mm,top=1.5mm,bottom=1.5mm,
  title={#1},fonttitle=\bfseries
}
\newtcolorbox{cautionbox}[1][]{
  enhanced,breakable,colback=softred,colframe=brick,
  boxrule=0.7pt,arc=1.5mm,left=2mm,right=2mm,top=1.5mm,bottom=1.5mm,
  title={#1},fonttitle=\bfseries
}

\lstdefinestyle{pythonstyle}{
  language=Python,
  basicstyle=\ttfamily\footnotesize,
  columns=fullflexible,
  frame=single,
  breaklines=true,
  showstringspaces=false,
  backgroundcolor=\color{softgray},
  rulecolor=\color{gray!55},
  keywordstyle=\bfseries,
  commentstyle=\itshape
}

\title{\vspace{-1.6cm}\textbf{From a Real Rank-One Relation to Positive-Density Torsion Transversality}\\[0.35em]
\Large A Kummer--Chebotarev and Algebraic \(K\)-Theory Progress Report\\
\large on the Alaoglu--Erd\H{o}s Integer-Exponent Conjecture}
\author{OpenAI GPT-5.6 Pro\\
\small New material beyond the attached \ReportA}
\date{22 August 2026}

\begin{document}
\maketitle

\begin{abstract}
The Alaoglu--Erd\H{o}s integer-exponent conjecture asks whether
\[
 2^x,3^x\in\Z \quad\Longrightarrow\quad x\in\Z
 \qquad (x\in\mathbb R).
\]
The attached \ReportA{} already develops an unusually broad corpus: exact conditional
structure, machine-checked primitive-output normalization, transcendence reductions,
finite-difference and determinant methods, \(p\)-adic diagnostics, interpolation,
symbolic dynamics, canonical products, and twenty-seven continuation reports.  The present
report deliberately does not reproduce those arguments.  It imports only two exact black
boxes from that corpus and develops a new finite-torsion branch.

For a putative least counterexample \(\beta\), write \(M=2^\beta\) and \(A=3^\beta\), and
form the homogeneous quadratic prime-exponent polynomial
\[
 Q_{M,A}(X)=L_M(X)X_3-L_A(X)X_2,
 \qquad L_u(X)=\sum_p v_p(u)X_p.
\]
The primitive-output theorem in \ReportA{} implies the new identity
\[
 \cont(Q_{M,A})=\gcd(d,e,a-b)=1
\]
for its simultaneous normalization \(M=2^aw^d\), \(A=3^bv^e\).  Thus \(Q_{M,A}\) remains
nonzero modulo every prime.  Kummer theory and Chebotarev then turn this elementary content
calculation into uniform local theorems at every prime torsion level.  At the quadratic level,
a multiquadratic Chebotarev argument and the minimum-distance bound for degree-two Boolean
polynomials force detection at density at least \(1/4\) among all rational primes.  For every
odd prime \(q\), a positive-density set of primes \(\ell\equiv1\pmod q\) detects rank two in
the \(q\)-th power-residue characters of \((2,3;M,A)\).  The relative detecting density is at
least \((1-1/q)^2\); for cubic characters it is at least \(4/9\) among
\(\ell\equiv1\pmod3\), hence at least \(2/9\) among all primes.  For all but finitely many
odd \(q\), the rank-three and rank-four branches have exact relative densities
\(1-1/q\) and \(1-1/q-(q-1)/q^3\), respectively.  Under GRH for fixed-degree Kummer
fields, quadratic and cubic detecting primes satisfy \(\ell\ll(\log(6MA))^2\).

These are unconditional structural advances, but they do not prove the conjecture.  Instead,
they isolate a sharp missing bridge: the real logarithmic rank-one relation would have to force
vanishing of the cubic determinant on more than \(5/9\) of the primes
\(\ell\equiv1\pmod3\), contradicting Chebotarev.  A second new branch, through Milnor
\(K_2(\Q)\) and tame symbols, gives an exact congruence sieve and also proves a parity
obstruction: ordinary odd-primary cup products are alternating, whereas the Alaoglu--Erd\H{o}s
period is symmetric.  This explains why a direct reciprocity proof cannot simply be read off
from the real determinant.

No full proof is claimed.  Every theorem, conditional theorem, computation, and open transfer
target is labelled separately.
\end{abstract}

\begingroup
\small
\tableofcontents
\endgroup
\newpage

\section{Executive verdict and claim ledger}

\begin{statusbox}[Current status]
No proof of the Alaoglu--Erd\H{o}s conjecture was obtained.  The main advance is a new,
fully rigorous Kummer--Chebotarev theorem that forces positive-density finite torsion
transversality for every hypothetical least counterexample.  The theorem is compatible with
the conjecture being false, so it is a reduction and a source of certificates, not a
contradiction by itself.
\end{statusbox}

The report's claims are divided into four levels.

\begin{center}
\begin{tabularx}{0.96\textwidth}{@{}>{\bfseries}p{0.18\textwidth}X@{}}
\toprule
Level & Meaning \\
\midrule
Proved here & Complete paper proof is supplied from standard stated inputs. \\
Inherited & Used as a black box from the attached \ReportA; not reproved here. \\
Conditional & The conclusion is exact under a named hypothesis such as GRH or a proposed transfer theorem. \\
Computational & Reproducible finite experiment; never used as proof of an infinite statement. \\
Open target & A precise statement whose proof would close a remaining gap. \\
\bottomrule
\end{tabularx}
\end{center}

\begin{resultbox}[Principal new conclusions]
Assume the conjecture fails, let \(\beta\) be its least nonintegral solution, and put
\(M=2^\beta\), \(A=3^\beta\).
\begin{enumerate}
\item The quadratic prime-exponent obstruction \(Q_{M,A}\) is primitive over \(\Z\).
\item The quadratic-character determinant is nonzero at a set of rational primes of density
at least \(1/4\).
\item For every odd prime \(q\), the \(q\)-th power-residue determinant is nonzero on a
positive-density set of primes.  The absolute density is at least \((q-1)/q^2\).
\item At \(q=3\), the detecting set has relative density at least \(4/9\) inside
\(\ell\equiv1\pmod3\), and absolute density at least \(2/9\).
\item Under GRH, quadratic and cubic detecting primes exist below effective constants times
\((\log(6MA))^2\); unconditionally one obtains fixed powers of \(6MA\).
\item A theorem forcing cubic determinant zero on a relative density greater than \(5/9\)
would prove the conjecture.
\item The ordinary \(K_2\)/cup-product route sees an alternating symbol with the wrong sign
symmetry.  It yields congruence restrictions but not the real period relation.
\end{enumerate}
\end{resultbox}

The source audit used the decompressed attachment
\nolinkurl{alaoglu-erdos-unified-report-A.tex.gz}, containing 20,925 lines.  Its decompressed
SHA-256 digest is
\begin{center}
\texttt{6bb33ec36c1d548ae8f343681bda054e752d62a24a0ac4157e4471dbd0a5c565}.
\end{center}
A structural and keyword audit found extensive treatment of real and \(p\)-adic logarithmic
determinants, but no Chebotarev density theorem, power-residue character determinant, Milnor
\(K_2\), tame-symbol, or norm-residue formulation.  Those are the delta developed below.

\section{Minimal inherited input from \ReportA}
\label{sec:inherited}

Only the following two pieces of the attachment are imported.

\subsection{Least counterexample and simultaneous primitive outputs}

Assume failure and let \(\beta>0\) be the least nonintegral real number with
\(2^\beta,3^\beta\in\Z\).  Set
\[
 M=2^\beta,\qquad A=3^\beta.
\]
The attachment proves that \(\beta\) is irrational and supplies factorizations
\begin{equation}
 M=2^a w^d,\qquad A=3^b v^e,\label{eq:normalization}
\end{equation}
where
\begin{itemize}
\item \(a,b\ge0\), \(d,e\ge1\);
\item \(w>1\) is odd and power-primitive;
\item \(v>1\) is power-primitive and \(3\nmid v\);
\item \(a=0\) or \(b=0\); and
\item the exact affine-primitivity identity holds:
\begin{equation}
 \gcd\bigl(d,e,a-b\bigr)=1.\label{eq:gcdone}
\end{equation}
\end{itemize}
Here power-primitive means that the gcd of the positive prime valuations is one.  Thus if
\[
 w=\prod_{p\mid w}p^{r_p},\qquad v=\prod_{p\mid v}p^{s_p},
\]
then \(\gcd_p r_p=\gcd_p s_p=1\).

\subsection{The depth-two symmetric tensor has rank three or four}

Let
\[
 \Gamma=\Q_{>0}^{\times}\cong\bigoplus_p\Z e_p,
 \qquad \nu(u)=\sum_pv_p(u)e_p.
\]
The real relation
\[
 \frac{\log M}{\log2}=\frac{\log A}{\log3}=\beta
\]
is equivalent to the vanishing of the period of
\begin{equation}
 \tau_{M,A}=\nu(M)\odot e_3-\nu(A)\odot e_2
 \in\Sym^2(\Gamma\otimes\Q),\label{eq:symtensor}
\end{equation}
under \(e_p\odot e_r\mapsto\log p\log r\).  The attachment combines a finite rank
classification with Baker's linear-independence theorem \cite{Baker1966} to show that, for a hypothetical
counterexample, \(\tau_{M,A}\) has quadratic rank three or four.  In the rank-four branch it
is the pullback of the split hyperbolic form \(z_1z_2-z_3z_4\).

Nothing else from the 27-report corpus is needed for the new theorems.

\section{The primitive quadratic obstruction}
\label{sec:primitiveQ}

Introduce one variable \(X_p\) for each prime and, for \(u\in\Q_{>0}^{\times}\), the finite
linear form
\[
 L_u(X)=\sum_p v_p(u)X_p.
\]
The polynomial attached to an output pair is
\begin{equation}
 Q_{M,A}(X)=L_M(X)X_3-L_A(X)X_2.\label{eq:Qdef}
\end{equation}
It is the diagonal evaluation of the symmetric tensor in \eqref{eq:symtensor}.  In
particular,
\begin{equation}
 Q_{M,A}\bigl((\log p)_p\bigr)
 =\log M\log3-\log A\log2.\label{eq:realperiod}
\end{equation}
Thus a counterexample gives a nonzero rational quadratic polynomial having one distinguished
transcendental zero.

For an integral polynomial in finitely many variables, \(\cont(Q)\) denotes the gcd of its
coefficients, taken nonnegative.

\begin{theorem}[Exact content formula]
\label{thm:content}
For the simultaneous normalization \eqref{eq:normalization},
\begin{equation}
 \cont(Q_{M,A})=\gcd(d,e,a-b).\label{eq:content}
\end{equation}
Consequently a least hypothetical counterexample has \(\cont(Q_{M,A})=1\).
\end{theorem}

\begin{proof}
Write
\[
 W(X)=\sum_{p\mid w}r_pX_p,
 \qquad
 V(X)=\sum_{p\mid v}s_pX_p.
\]
Since \(w\) is odd, \(W\) has no \(X_2\)-term; since \(3\nmid v\), \(V\) has no
\(X_3\)-term.  From \eqref{eq:normalization},
\begin{align*}
 L_M&=aX_2+dW,\\
 L_A&=bX_3+eV,
\end{align*}
and hence
\begin{equation}
 Q_{M,A}=(a-b)X_2X_3+dW(X)X_3-eV(X)X_2.\label{eq:Qexpanded}
\end{equation}
There is no hidden cancellation in this expansion.  The first displayed monomial has
coefficient \(a-b\).  The coefficients contributed by \(dWX_3\) have gcd \(d\), because
\(\gcd_pr_p=1\).  The coefficients contributed by \(eVX_2\) have gcd \(e\), because
\(\gcd_ps_p=1\).  Even if a prime divides both \(w\) and \(v\), the corresponding monomials
are \(X_pX_3\) and \(X_pX_2\), so they remain distinct.  If \(3\mid w\) or \(2\mid v\), the
extra monomials are \(X_3^2\) or \(X_2^2\), again distinct.  Therefore the gcd of all
coefficients is precisely \(\gcd(d,e,a-b)\).  Equation \eqref{eq:gcdone} proves the final
claim.
\end{proof}

\begin{corollary}[Universal nonvanishing after reduction]
\label{cor:modqnonzero}
For a least hypothetical counterexample, the reduction of \(Q_{M,A}\) is a nonzero quadratic
polynomial over \(\F_q\) for every prime \(q\).
\end{corollary}

\begin{proof}
A primitive integral polynomial cannot have every coefficient divisible by \(q\).
\end{proof}

\begin{remark}
The gcd-one output theorem in \ReportA{} was originally a descent statement.  The new point
is that it is exactly the coefficient-content statement needed to keep the symmetric
obstruction alive at every torsion level.  This is the hinge between the real-period problem
and Kummer theory.
\end{remark}

\subsection{Standard pairs are exactly the zero-polynomial pairs}

\begin{proposition}[Formal standardness criterion]
\label{prop:Qzero}
For positive integers \(M,A\), the following are equivalent:
\begin{enumerate}
\item \(Q_{M,A}=0\) over \(\Z\);
\item there is an integer \(n\ge0\) such that \(M=2^n\) and \(A=3^n\).
\end{enumerate}
\end{proposition}

\begin{proof}
The reverse implication is immediate.  Conversely, compare coefficients in
\(L_MX_3-L_AX_2\).  The coefficient of \(X_pX_3\), for \(p\notin\{2,3\}\), is
\(v_p(M)\), while the coefficient of \(X_pX_2\) is \(-v_p(A)\); hence both vanish.  The
coefficients of \(X_3^2\) and \(X_2^2\) give \(v_3(M)=v_2(A)=0\).  Finally the coefficient
of \(X_2X_3\) gives \(v_2(M)=v_3(A)=n\).
\end{proof}

The contrast is already sharp:
\[
 \begin{array}{c|c}
 \text{integer-exponent output pair} & Q_{M,A}=0,\\
 \text{least hypothetical noninteger output pair} & \cont(Q_{M,A})=1.
 \end{array}
\]
The next sections make this integral dichotomy visible in reductions modulo rational primes.

\section{Kummer characters and the finite determinant}
\label{sec:kummer}

Fix an odd prime \(q\).  Let
\[
 \overline\Gamma_q=\Q_{>0}^{\times}/(\Q_{>0}^{\times})^q
 \cong\bigoplus_p\F_q e_p
\]
and let
\begin{equation}
 W_q=\Span_{\F_q}\bigl\{[2],[3],[M],[A]\bigr\}
 \subset\overline\Gamma_q,
 \qquad r_q=\dim_{\F_q}W_q.\label{eq:Wq}
\end{equation}
Because the classes of \(2\) and \(3\) are independent, \(2\le r_q\le4\).

Put \(K_q=\Q(\zeta_q)\), where \(\zeta_q\) is a primitive \(q\)-th root of unity, and
consider
\begin{equation}
 \mathcal L_q
 =K_q\bigl(2^{1/q},3^{1/q},M^{1/q},A^{1/q}\bigr).\label{eq:Lq}
\end{equation}

\begin{lemma}[No loss of rational Kummer rank]
\label{lem:injection}
The natural map
\[
 \Q^{\times}/(\Q^{\times})^q\longrightarrow K_q^{\times}/(K_q^{\times})^q
\]
is injective.  In particular,
\begin{equation}
 \Gal(\mathcal L_q/K_q)\cong\operatorname{Hom}_{\F_q}(W_q,\mu_q)
 \quad\text{and}\quad
 [\mathcal L_q:K_q]=q^{r_q}.\label{eq:kummergal}
\end{equation}
\end{lemma}

\begin{proof}
Suppose \(u\in\Q^\times\) becomes a \(q\)-th power in \(K_q\), say \(u=y^q\).  Taking
norms to \(\Q\) gives
\[
 u^{q-1}=N_{K_q/\Q}(y)^q.
\]
Thus the class of \(u\) in \(\Q^\times/(\Q^\times)^q\) is annihilated by \(q-1\).  It is
also \(q\)-torsion, and \(\gcd(q,q-1)=1\), so the class is zero.  The Kummer description
then follows from the standard Kummer theorem over a field containing \(\mu_q\)
\cite{MilneFT}.
\end{proof}

Let \(\ell\) be a rational prime satisfying
\[
 \ell\equiv1\pmod q,
 \qquad \ell\nmid 6MAq.
\]
Choose \(\omega_\ell\in\F_\ell^\times\) of exact order \(q\).  For
\(u\in\Q^\times\) with \(\ell\nmid u\), define the additive \(q\)-th power-residue
character \(\chi_{\ell,q}(u)\in\F_q\) by
\begin{equation}
 u^{(\ell-1)/q}\equiv
 \omega_\ell^{\,\chi_{\ell,q}(u)}\pmod\ell.\label{eq:qchar}
\end{equation}
It is a homomorphism \(\overline\Gamma_q\to\F_q\).  Replacing \(\omega_\ell\) by another
generator scales every character value by one common nonzero scalar.

\begin{definition}[Power-residue determinant]
\label{def:delta}
Set
\begin{equation}
 \Delta_{\ell,q}(M,A)
 =\chi_{\ell,q}(M)\chi_{\ell,q}(3)
  -\chi_{\ell,q}(A)\chi_{\ell,q}(2)
 \in\F_q.\label{eq:delta}
\end{equation}
Its vanishing is independent of the choice of \(\omega_\ell\), because a change of generator
multiplies \(\Delta_{\ell,q}\) by a nonzero square.
\end{definition}

\begin{proposition}[The determinant is the quadratic obstruction evaluated at Frobenius]
\label{prop:deltaQ}
For every eligible \(\ell\),
\begin{equation}
 \Delta_{\ell,q}(M,A)=Q_{M,A}(\chi_{\ell,q}).\label{eq:deltaQ}
\end{equation}
Moreover, up to the common scalar ambiguity just described, \(\chi_{\ell,q}|_{W_q}\) is the
Kummer character of Frobenius at a prime of \(K_q\) above \(\ell\).
\end{proposition}

\begin{proof}
The first identity follows by substituting \(X_p=\chi_{\ell,q}(p)\) into
\eqref{eq:Qdef} and using additivity on prime valuations.  The second statement is the usual
power-residue description of Frobenius in a Kummer extension.
\end{proof}

For a standard pair \((M,A)=(2^n,3^n)\), both rows of the character matrix are proportional:
\[
 \begin{pmatrix}
 \chi(2)&\chi(M)\\
 \chi(3)&\chi(A)
 \end{pmatrix}
 =
 \begin{pmatrix}
 \chi(2)&n\chi(2)\\
 \chi(3)&n\chi(3)
 \end{pmatrix},
\]
so \(\Delta_{\ell,q}=0\) for every \(q\) and every eligible \(\ell\).  A least hypothetical
counterexample will exhibit the opposite behavior on a positive-density set at every odd
level.

\section{Chebotarev forces positive-density torsion transversality}
\label{sec:chebotarev}

The extension \(\mathcal L_q/\Q\) is Galois.  Its Galois group is a semidirect product
\begin{equation}
 \Gal(\mathcal L_q/\Q)
 \cong W_q^\vee\rtimes\F_q^\times,
 \qquad W_q^\vee=\operatorname{Hom}_{\F_q}(W_q,\F_q),\label{eq:semidirect}
\end{equation}
where \(\F_q^\times\) scales the Kummer character.  Primes
\(\ell\equiv1\pmod q\) have trivial cyclotomic component, and the property
\(Q_{M,A}(\chi)\ne0\) is stable under scaling \(\chi\mapsto c\chi\).

\begin{theorem}[Exact Chebotarev density formula]
\label{thm:densityexact}
Suppose the reduction of \(Q_{M,A}\) modulo \(q\) is nonzero on \(W_q^\vee\).  Then the set
\[
 \mathscr D_q(M,A)=
 \bigl\{\ell:\ell\equiv1\pmod q,\ \ell\nmid6MAq,
              \ \Delta_{\ell,q}(M,A)\ne0\bigr\}
\]
has natural density
\begin{equation}
 \dens\mathscr D_q(M,A)
 =\frac{\#\{\chi\in W_q^\vee:Q_{M,A}(\chi)\ne0\}}
        {(q-1)q^{r_q}}.\label{eq:absdensity}
\end{equation}
Relative to the primes \(\ell\equiv1\pmod q\), its density is
\begin{equation}
 \dens_{\ell\equiv1(q)}\mathscr D_q(M,A)
 =\frac{\#\{\chi\in W_q^\vee:Q_{M,A}(\chi)\ne0\}}
        {q^{r_q}}.\label{eq:reldensity}
\end{equation}
\end{theorem}

\begin{proof}
Let
\(S_q=\{\chi\in W_q^\vee:Q_{M,A}(\chi)\ne0\}\).
Because \(Q(c\chi)=c^2Q(\chi)\), \(S_q\) is stable under the conjugation action of
\(\F_q^\times\) in \eqref{eq:semidirect}; it is therefore a union of conjugacy classes in
\(\Gal(\mathcal L_q/\Q)\).  Chebotarev's density theorem \cite{MilneANT,SerreCheb} assigns it density
\(\#S_q/\#\Gal(\mathcal L_q/\Q)=\#S_q/((q-1)q^{r_q})\).  The cyclotomic splitting
condition \(\ell\equiv1\pmod q\) itself has density \(1/(q-1)\), giving
\eqref{eq:reldensity}.
\end{proof}

The required finite-field estimate is sharp in the worst case.

\begin{lemma}[Universal zero bound for a nonzero quadratic form]
\label{lem:quadzeros}
Let \(q\) be odd, let \(r\ge2\), and let \(Q\) be a nonzero homogeneous quadratic
polynomial on \(\F_q^r\).  Then
\begin{equation}
 \#\{x\in\F_q^r:Q(x)=0\}\le(2q-1)q^{r-2}.\label{eq:quadzerobound}
\end{equation}
Equivalently,
\begin{equation}
 \frac{\#\{x:Q(x)\ne0\}}{q^r}
 \ge\left(1-\frac1q\right)^2.\label{eq:quadnonzero}
\end{equation}
Equality in \eqref{eq:quadzerobound} occurs for a split rank-two form such as \(x_1x_2\).
\end{lemma}

\begin{proof}
Pass to the nondegenerate quotient by the radical and let \(\rho\) be the rank.  For
\(\rho=1\), the zero set has \(q^{r-1}\) points.  For \(\rho=2\), an anisotropic form has
\(q^{r-2}\) zeros, while a split form has
\((2q-1)q^{r-2}\) zeros.  For odd \(\rho\ge3\), the standard count is \(q^{r-1}\).  For
even \(\rho=2m\ge4\), it is
\[
 q^{r-1}+\varepsilon(q-1)q^{r-m-1},\qquad \varepsilon\in\{-1,1\},
\]
which is no larger than the split rank-two count.  This proves the bound.  Subtracting from
\(q^r\) yields \eqref{eq:quadnonzero}.
\end{proof}

\begin{theorem}[Universal torsion transversality]
\label{thm:universal}
Let \((M,A)\) arise from a least hypothetical counterexample.  For every odd prime \(q\),
\begin{align}
 \dens_{\ell\equiv1(q)}\mathscr D_q(M,A)
 &\ge\left(1-\frac1q\right)^2,\label{eq:relativebound}\\
 \dens\mathscr D_q(M,A)
 &\ge\frac{q-1}{q^2}.\label{eq:absolutebound}
\end{align}
In particular, for cubic residue characters,
\begin{align}
 \dens_{\ell\equiv1(3)}\mathscr D_3(M,A)&\ge\frac49,\label{eq:cubicrelative}\\
 \dens\mathscr D_3(M,A)&\ge\frac29.\label{eq:cubicabsolute}
\end{align}
\end{theorem}

\begin{proof}
By \Cref{cor:modqnonzero}, the integral polynomial remains nonzero modulo \(q\).  It belongs
to \(\Sym^2(W_q)\), and the inclusion \(W_q\hookrightarrow\overline\Gamma_q\) induces an
injection on symmetric squares.  Because \(q\) is odd, a nonzero symmetric tensor defines a
nonzero quadratic polynomial on \(W_q^\vee\).  Apply \Cref{lem:quadzeros} to
\eqref{eq:reldensity}; division by \(q-1\) gives the absolute density.
\end{proof}

\begin{resultbox}[Cubic Chebotarev certificate]
A least counterexample would be locally rank two in cubic power-residue characters at no less
than \(2/9\) of all rational primes.  This lower bound is unconditional and uniform in the
size and factorization of \(M\) and \(A\).
\end{resultbox}

\subsection{A local characterization of standard output pairs}

The preceding argument applies to any positive integer pair, with the coefficient content
measuring the exceptional torsion levels.

\begin{theorem}[Torsion-profile dichotomy]
\label{thm:profile}
Let \(M,A\in\N\), and suppose \(Q_{M,A}\ne0\).  Put \(c=\cont(Q_{M,A})\).  For every odd
prime \(q\):
\begin{enumerate}
\item if \(q\mid c\), then \(\Delta_{\ell,q}(M,A)=0\) for every eligible \(\ell\);
\item if \(q\nmid c\), then the detecting set has the positive density in
\eqref{eq:absdensity}, and in particular satisfies \eqref{eq:absolutebound}.
\end{enumerate}
Consequently the following are equivalent:
\begin{enumerate}[label=(\alph*)]
\item \(M=2^n\) and \(A=3^n\) for some \(n\ge0\);
\item the determinant vanishes at every eligible prime for every odd \(q\);
\item there are infinitely many odd \(q\) for which the determinant vanishes on a density-one
subset of the primes \(\ell\equiv1\pmod q\).
\end{enumerate}
\end{theorem}

\begin{proof}
Reduction modulo \(q\) is zero exactly when \(q\) divides every coefficient, which is exactly
\(q\mid c\).  Otherwise \Cref{thm:densityexact,lem:quadzeros} apply.  If (c) holds for a
nonzero \(Q\), infinitely many primes would divide its nonzero content, impossible.  Hence
\(Q=0\), and \Cref{prop:Qzero} gives (a).  The other implications are immediate.
\end{proof}

\begin{remark}
This is a genuine local characterization of the integer-exponent ray.  It does not use the
real equation \eqref{eq:realperiod}; rather, it says that a nonstandard integer pair cannot
hide from all odd torsion characters.  The Alaoglu--Erd\H{o}s difficulty is to make the real
rank-one relation interact with this local visibility.
\end{remark}

\section{Rank-stratified exact densities}
\label{sec:rankdensity}

The universal lower bound deliberately allows every possible rank drop modulo \(q\).  The
rank theorem inherited from \ReportA{} gives more for all but finitely many \(q\).

Let \(\rho\in\{3,4\}\) be the rational rank of the symmetric tensor
\(\tau_{M,A}\).  Every nonzero minor witnessing this rank is an integer, so only finitely many
primes divide all such minors.  Away from those primes, reduction preserves the rank and the
multiplicative span.

\begin{theorem}[Exact generic torsion densities]
\label{thm:genericdensity}
For all but finitely many odd primes \(q\), the following alternatives hold.
\begin{enumerate}
\item If \(\rho=3\), then \(r_q=3\) and
\begin{equation}
 \dens_{\ell\equiv1(q)}\mathscr D_q(M,A)=1-\frac1q,
 \qquad
 \dens\mathscr D_q(M,A)=\frac1q.\label{eq:rank3density}
\end{equation}
\item If \(\rho=4\), then \(r_q=4\), the reduced form is split hyperbolic, and
\begin{align}
 \dens_{\ell\equiv1(q)}\mathscr D_q(M,A)
 &=1-\frac1q-\frac{q-1}{q^3},\label{eq:rank4relative}\\
 \dens\mathscr D_q(M,A)
 &=\frac{q^3-q^2-q+1}{(q-1)q^3}.\label{eq:rank4absolute}
\end{align}
\end{enumerate}
In either branch, the relative detecting density tends to one as \(q\to\infty\).
\end{theorem}

\begin{proof}
For the rank-three branch, the finite quadratic form has a nondegenerate ternary quotient.
Every nondegenerate quadratic form in an odd number of variables has exactly \(q^{r_q-1}\)
zeros, hence nonzero proportion \(1-1/q\).

For the rank-four branch, the four defining linear forms remain independent and
\(Q\) is equivalent to \(z_1z_2-z_3z_4\).  A split nondegenerate four-variable form has
\[
 q^3+(q-1)q
\]
zeros.  Division by \(q^4\), followed by \Cref{thm:densityexact}, gives
\eqref{eq:rank4relative} and \eqref{eq:rank4absolute}.
\end{proof}

At a nonexceptional cubic level this gives the more precise table
\begin{center}
\begin{tabular}{@{}lcc@{}}
\toprule
Generic rational branch & relative density in \(\ell\equiv1\pmod3\) & absolute density \\
\midrule
rank three & \(2/3\) & \(1/3\) \\
rank four & \(16/27\) & \(8/27\) \\
\bottomrule
\end{tabular}
\end{center}
The robust theorem \eqref{eq:cubicabsolute} uses only \(2/9\), because \(q=3\) may be one of
the finitely many rank-drop primes.

\begin{remark}[Why taking \(q\) large does not itself prove anything]
The relative detecting density approaches one, but the ambient congruence class
\(\ell\equiv1\pmod q\) has density \(1/(q-1)\).  The absolute density is therefore of order
\(1/q\).  No contradiction follows by letting \(q\) grow; one needs a theorem that transfers
the real period zero into systematic finite determinant zeros.
\end{remark}

\section{Characteristic two: Boolean Kummer transversality}
\label{sec:char2}

The prime two is not covered by the odd-character quadratic-form count, but it gives a
slightly stronger absolute density.  Put
\[
 W_2=\Span_{\F_2}\{[2],[3],[M],[A]\}
 \subset \Q_{>0}^{\times}/(\Q_{>0}^{\times})^2,
 \qquad r_2=\dim_{\F_2}W_2,
\]
and
\begin{equation}
 \mathcal L_2=\Q(\sqrt2,\sqrt3,\sqrt M,\sqrt A).
 \label{eq:L2}
\end{equation}
Then \(\Gal(\mathcal L_2/\Q)\cong W_2^\vee\) and
\([\mathcal L_2:\Q]=2^{r_2}\le16\).

For an odd prime \(\ell\nmid6MA\), define the additive quadratic character by
\[
 \left(\frac{u}{\ell}\right)=(-1)^{\chi_{\ell,2}(u)},
 \qquad \chi_{\ell,2}(u)\in\F_2,
\]
and set
\begin{equation}
 \Delta_{\ell,2}(M,A)
 =\chi_{\ell,2}(M)\chi_{\ell,2}(3)
  +\chi_{\ell,2}(A)\chi_{\ell,2}(2).
 \label{eq:delta2}
\end{equation}
The plus sign is the reduction of the minus sign in \eqref{eq:delta}.  As before,
\(\Delta_{\ell,2}=Q_{M,A}(\chi_{\ell,2})\), and Chebotarev gives
\begin{equation}
 \dens\{\ell:\ell\nmid6MA,\ \Delta_{\ell,2}\ne0\}
 =\frac{\#\{\chi\in W_2^\vee:Q_{M,A}(\chi)\ne0\}}{2^{r_2}}.
 \label{eq:density2exact}
\end{equation}

\begin{lemma}[Boolean degree-two support bound]
\label{lem:booleanweight}
Let \(P\) be a nonzero homogeneous quadratic polynomial over \(\F_2\) on an
\(r\)-dimensional vector space, with \(r\ge2\).  Then
\begin{equation}
 \#\{x\in\F_2^r:P(x)\ne0\}\ge2^{r-2}.
 \label{eq:booleanweight}
\end{equation}
The bound is sharp, for example for \(P=x_1x_2\).
\end{lemma}

\begin{proof}
Choose coordinates and replace every \(x_i^2\) by \(x_i\).  This gives the unique
multilinear polynomial defining the same function on the Boolean cube.  It is still nonzero:
the square coefficients become distinct linear coefficients, while the mixed coefficients
remain distinct quadratic coefficients.

More generally, a nonzero multilinear polynomial of degree at most \(d\) on
\(\F_2^r\) has Hamming weight at least \(2^{r-d}\).  Induct on \(r\), writing
\(P=P_0+x_rP_1\).  If \(P_1=0\), the weight doubles from dimension \(r-1\).  If
\(P_1\ne0\), then coordinatewise
\(\operatorname{wt}(P_0)+\operatorname{wt}(P_0+P_1)
 \ge\operatorname{wt}(P_1)\), and induction at degree \(d-1\) gives the same lower
bound.  Taking \(d=2\) proves \eqref{eq:booleanweight}.
\end{proof}

\begin{theorem}[Quadratic-character transversality]
\label{thm:quadratictransversality}
For a least hypothetical counterexample,
\begin{equation}
 \dens\{\ell:\ell\nmid6MA,\ \Delta_{\ell,2}(M,A)\ne0\}
 \ge\frac14.
 \label{eq:quadraticdensity}
\end{equation}
Thus at least one quarter of all rational primes detect nonstandardness through the four
Legendre-symbol bits of \((2,3,M,A)\).
\end{theorem}

\begin{proof}
Primitivity of \(Q_{M,A}\) keeps its reduction modulo two nonzero.  The tensor lies in
\(\Sym^2(W_2)\); hence its restriction is a nonzero homogeneous quadratic polynomial on
\(W_2^\vee\).  In characteristic two a nonzero homogeneous quadratic remains a nonzero
function on the Boolean cube by the first paragraph of \Cref{lem:booleanweight}.  Apply
\eqref{eq:density2exact} and \Cref{lem:booleanweight}.
\end{proof}

\begin{resultbox}[Strongest uniform absolute-density certificate]
A least counterexample would be detected by quadratic characters at density at least
\(1/4\) of all primes, and by cubic characters at density at least \(2/9\).  The quadratic
level gives the larger absolute density and a smaller field; the cubic level gives the sharper
transfer threshold, because its zero density is at most \(5/9\) inside one congruence class.
\end{resultbox}

\section{Effective detecting primes}
\label{sec:effective}

The quadratic and cubic levels both have fixed degree.  Along with \eqref{eq:L2}, set
\begin{equation}
 \mathcal L_3=\Q\bigl(\zeta_3,2^{1/3},3^{1/3},M^{1/3},A^{1/3}\bigr).
 \label{eq:L3}
\end{equation}
Then
\begin{equation}
 [\mathcal L_2:\Q]\le16,
 \qquad
 [\mathcal L_3:\Q]\le 2\cdot3^4=162.
 \label{eq:degrees}
\end{equation}
Only primes dividing \(6MA\) can ramify.  The discriminants of the defining radical
polynomials and the standard compositum inequality \cite{MilneANT,Neukirch} therefore give
absolute effective constants \(C_{0,2},C_{0,3}\) such that
\begin{equation}
 \log\abs{D_{\mathcal L_j}}\le C_{0,j}\log(6MA),
 \qquad j\in\{2,3\}.
 \label{eq:discbounds}
\end{equation}
The constants are independent of the number of prime factors and of the exponent vectors of
\(M\) and \(A\).

\begin{theorem}[Effective low-torsion witnesses]
\label{thm:effectivewitness}
For a least hypothetical counterexample there are eligible detecting primes
\(\ell_2,\ell_3\) with
\[
 \Delta_{\ell_2,2}(M,A)\ne0,
 \qquad
 \Delta_{\ell_3,3}(M,A)\ne0.
\]
More precisely, for each \(j\in\{2,3\}\):
\begin{enumerate}
\item assuming GRH for the Dedekind zeta function of \(\mathcal L_j\),
\begin{equation}
 \ell_j\ll \bigl(\log(6MA)\bigr)^2,
 \label{eq:GRHbound}
\end{equation}
with an absolute effective implied constant;
\item unconditionally,
\begin{equation}
 \ell_j\le(6MA)^{C_j}
 \label{eq:uncondbound}
\end{equation}
for an absolute effective constant \(C_j\).
\end{enumerate}
Here \(\ell_2\nmid6MA\), while \(\ell_3\equiv1\pmod3\) and \(\ell_3\nmid6MA\).
\end{theorem}

\begin{proof}
At level two, \Cref{thm:quadratictransversality} supplies a nonempty subset of
\(\Gal(\mathcal L_2/\Q)\) on which \(Q\ne0\).  At level three,
\Cref{thm:universal} supplies a Kummer character \(\chi\in W_3^\vee\) with
\(Q_{M,A}(\chi)\ne0\); its orbit under \(\F_3^\times\) is a nonempty conjugacy-stable subset
of \(\Gal(\mathcal L_3/\Q)\).  Effective Chebotarev
\cite{LagariasOdlyzko,GrenieMolteni,ThornerZaman,KadiriWong} applies to either subset.  In the
cubic case its cyclotomic component is trivial, so the selected prime is
\(1\pmod3\).

Let \(S_0\) be the finite set of primes dividing \(6MA\).  Use the standard finite-exception
variant of effective Chebotarev.  In a weighted prime count, deleting \(S_0\) costs at most
\(\sum_{p\in S_0}\log p\le\log(6MA)\).  Under GRH, a first prime outside \(S_0\) in the
chosen nonempty conjugacy-stable subset is therefore bounded by a constant times
\[
 \bigl(\log\abs{D_{\mathcal L_j}}+[\mathcal L_j:\Q]+\log(6MA)\bigr)^2.
\]
Insert \eqref{eq:degrees}--\eqref{eq:discbounds}.  The unconditional effective counting
theorem, with the same finite subtraction, gives a fixed power of
\(\abs{D_{\mathcal L_j}}\,6MA\); another use of \eqref{eq:discbounds} gives
\eqref{eq:uncondbound}.  The selected primes lie outside \(S_0\), so they are eligible.
\end{proof}

\begin{resultbox}[Interpretation]
Under GRH, any explicit least counterexample would carry two polylogarithmic-size local
certificates that it is not on the integer-exponent ray.  The quadratic certificate needs only
four Legendre symbols in a field of degree at most 16.  The cubic certificate needs one cubic
root of unity modulo \(\ell\) and four exponentiations in a field of degree at most 162.
Neither certificate contradicts the pair's real logarithmic relation.
\end{resultbox}

\begin{remark}[Uniformity in support]
A naive radical extension adjoining a square or cube root of every prime factor of \(MA\)
would have degree growing exponentially with the support.  The fields \(\mathcal L_2\) and
\(\mathcal L_3\) adjoin only the four radicands \(2,3,M,A\); their degrees are uniformly
bounded by 16 and 162.  This fixed-degree compression is essential for \eqref{eq:GRHbound}.
\end{remark}

\section{The exact missing bridge: from the real place to cubic density}
\label{sec:bridge}

The Chebotarev theorem determines the finite distribution from the algebraic classes of
\(2,3,M,A\).  The real counterexample equation is
\begin{equation}
 \log M\log3-\log A\log2=0.\label{eq:realzeroagain}
\end{equation}
There is no formal reduction of the positive real logarithms modulo \(q\), and global
reciprocity does not identify the real magnitude character with a finite Kummer character.
A proof must therefore supply a genuinely new transfer theorem.

\begin{criterion}[Cubic density-threshold criterion]
\label{crit:threshold}
The Alaoglu--Erd\H{o}s conjecture follows from the following statement:

\begin{quote}
Whenever positive integers \(M,A\) satisfy \eqref{eq:realzeroagain} and arise as the least
nonintegral output pair, the determinant \(\Delta_{\ell,3}(M,A)\) vanishes on a subset of the
primes \(\ell\equiv1\pmod3\) having relative lower density strictly greater than \(5/9\).
\end{quote}
\end{criterion}

\begin{proof}
\Cref{thm:universal} says that the complementary nonvanishing set has relative density at
least \(4/9\).  Hence the zero set has relative density at most \(5/9\), contradicting the
proposed transfer statement.
\end{proof}

The threshold is robust: it uses no rank preservation modulo three.  In a nonexceptional
rank-three branch the permitted zero density is exactly \(1/3\); in a nonexceptional
rank-four branch it is exactly \(11/27\).  Thus additional control of the cubic rank could
lower the required transfer threshold.

\begin{gapbox}[Core open problem isolated by this report]
Prove any mechanism by which the archimedean period zero \eqref{eq:realzeroagain} forces a
statistical excess of finite Kummer rank drops.  A density greater than \(5/9\) at the cubic
level is already sufficient.  A theorem forcing \(\Delta_{\ell,3}=0\) for all but finitely
many eligible primes would be far stronger than necessary.
\end{gapbox}

\subsection{Why ordinary continuity cannot supply the transfer}

The map \(u\mapsto\log u\) at the real place is continuous in the usual topology.  The map
\(u\mapsto\chi_{\ell,q}(u)\) is a torsion quotient of a finite residue group.  Real closeness
of two positive rationals says nothing about equality of their power-residue characters.
Indeed, the computations in \Cref{sec:computation} exhibit pairs whose real determinant is
smaller than \(10^{-11}\) while their cubic characters have the generic rank-four density.

Nor is there a missing application of the product formula.  The product formula relates
absolute values of one algebraic number over all places.  Here the real equation is a
quadratic relation among four logarithms, whereas \(\Delta_{\ell,q}\) is a quadratic
function of four independent Kummer characters.  Turning one into the other requires an
auxiliary algebraic number whose local behavior encodes the symmetric tensor, and no such
construction is currently known.

\section{A Chebotarev sieve for synchronized integer gaps}
\label{sec:gapsieve}

The attachment already studies convergents to a hypothetical \(\beta\) and the synchronized
integer gaps.  The new torsion theorem adds a precise forbidden-prime statement.

Let \(r/s\) be a rational approximation to \(\beta\), with \(s>0\), and define
\begin{equation}
 D_2(r,s)=M^s-2^r,
 \qquad
 D_3(r,s)=A^s-3^r.\label{eq:gaps}
\end{equation}
For a nonintegral \(\beta\), neither gap is zero when \(r/s\ne\beta\).

\begin{proposition}[Common gap primes lie on the Kummer quadric]
\label{prop:gapprime}
Let \(q\) be prime.  For \(q=2\), let \(\ell\nmid6MA\) be odd; for odd \(q\), let
\(\ell\) be eligible at level \(q\).  Suppose
\[
 \ell\mid D_2(r,s),\qquad \ell\mid D_3(r,s),
 \qquad q\nmid s.
\]
Then
\begin{equation}
 \Delta_{\ell,q}(M,A)=0.\label{eq:gapdelta0}
\end{equation}
\end{proposition}

\begin{proof}
Apply \(\chi_{\ell,q}\) to the two congruences
\(M^s\equiv2^r\pmod\ell\) and \(A^s\equiv3^r\pmod\ell\):
\[
 s\chi(M)=r\chi(2),
 \qquad
 s\chi(A)=r\chi(3).
\]
Taking the determinant gives \(s\Delta_{\ell,q}=0\).  Since \(s\ne0\) in \(\F_q\), the
claim follows.
\end{proof}

\begin{corollary}[Quadratic and cubic forbidden-prime sets]
\label{cor:forbidden}
Let
\[
 G(r,s)=\gcd\bigl(D_2(r,s),D_3(r,s)\bigr).
\]
If \(s\) is odd, then \(G(r,s)\) has no eligible prime divisor in the quadratic detecting
set from \Cref{thm:quadratictransversality}; for a least counterexample this excludes a set
of density at least \(1/4\).  If \(s\not\equiv0\pmod3\), then \(G(r,s)\) has no
prime divisor in \(\mathscr D_3(M,A)\), excluding a set of density at least \(2/9\).
\end{corollary}

Continued-fraction denominators have \(\gcd(s_n,s_{n+1})=1\), so infinitely many are odd
and infinitely many are not divisible by three.  Thus the two exclusions apply along infinite
subsequences of the best real approximants.

\begin{cautionbox}[Why the sieve still stops short]
An individual integer can be arbitrarily large while all its prime factors lie in the
complement of a fixed positive-density set.  Consequently the exclusion of
the quadratic or cubic detecting sets does not upper-bound \(G(r,s)\) strongly enough to repair the
exponential denominator loss diagnosed in \ReportA.  A successful argument would need an
additional distribution theorem for prime divisors of the synchronized gcd sequence, or a
lower bound forcing that gcd to sample the detecting Chebotarev set.
\end{cautionbox}

The proposition nevertheless sharpens the old statement ``real closeness does not impose a
common prime divisor.''  It identifies exactly which Frobenius classes a common divisor is
allowed to occupy.  This suggests a concrete analytic program: study the sequence
\(G(r_n,s_n)\) with a Chebotarev-weighted large sieve, rather than estimate only its size.

\section{Finite computational stress test}
\label{sec:computation}

This section is computational evidence only.  Its purpose is to test the formulas and to
measure how unrelated extreme real closeness can be from cubic torsion rank.

Let
\[
 \theta=\frac{\log3}{\log2}.
\]
For each \(2\le M\le12000\), the search examined the nearest integers to \(M^\theta\),
removed exact standard pairs, retained pairs whose \(Q_{M,A}\) has content one, and sorted by
\begin{equation}
 \delta_\infty(M,A)
 =\left|\frac{\log M}{\log2}-\frac{\log A}{\log3}\right|.\label{eq:realmismatch}
\end{equation}
For each retained pair, it found the first eligible cubic detecting prime and measured the
nonvanishing frequency among primes \(\ell\le200000\), \(\ell\equiv1\pmod3\).

\begin{experiment}[Near-miss pairs]
\label{comp:table}
The following six primitive pairs have very small real mismatch.  The last column is the
empirical cubic detection rate; all six have the generic split rank-four prediction
\(16/27=0.592592\ldots\).
\end{experiment}

\begin{center}
\small
\begin{tabular}{@{}rrrrr@{}}
\toprule
\(M\) & \(A\) & \(\delta_\infty(M,A)\) & first detecting \(\ell\) & empirical rate \\
\midrule
 9329 & 1959024 & \(7.6827\times10^{-12}\) & 13 & \(5399/8986=0.600824\) \\
 8427 & 1667418 & \(1.1870\times10^{-11}\) &  7 & \(5344/8988=0.594571\) \\
10855 & 2490703 & \(3.1786\times10^{-11}\) &  7 & \(5305/8987=0.590297\) \\
 9566 & 2038489 & \(3.4305\times10^{-11}\) & 19 & \(5293/8985=0.589093\) \\
10506 & 2364980 & \(4.2679\times10^{-11}\) &  7 & \(5326/8986=0.592700\) \\
 9517 & 2021964 & \(7.3365\times10^{-11}\) & 13 & \(5380/8985=0.598776\) \\
\bottomrule
\end{tabular}
\end{center}

For example,
\[
 9329=19\cdot491,
 \qquad
 1959024=2^4\cdot3\cdot40813.
\]
At \(\ell=13\), one choice of cubic character gives
\[
 \bigl(\chi(2),\chi(3),\chi(M),\chi(A)\bigr)=(1,1,0,1),
\]
so \(\Delta_{13,3}=0\cdot1-1\cdot1=2\ne0\) in \(\F_3\).

\begin{statusbox}[What the computation establishes]
The exact finite-field counts match the Chebotarev prediction, and the empirical prime
frequencies are close to \(16/27\).  The experiment also falsifies a tempting continuity
heuristic: a real determinant of size \(10^{-11}\) need not produce even a mild bias toward
cubic rank drop.  It says nothing about whether an exact real zero can occur.
\end{statusbox}

\section{Milnor \texorpdfstring{\(K_2\)}{K2}, tame symbols, and the parity barrier}
\label{sec:K2}

The Kummer determinant invites a reciprocity-theoretic attack.  This section carries that
idea to its exact endpoint.  It yields a valid congruence sieve but also exposes a structural
sign obstruction.

\subsection{Symmetric versus alternating degree two}

In additive prime-exponent notation, the target tensor is
\[
 \tau_{\mathrm{sym}}=\nu(M)\odot e_3-\nu(A)\odot e_2.
\]
If one keeps the displayed signs and antisymmetrizes, one obtains
\begin{equation}
 \tau^-_{\wedge}=\nu(M)\wedge e_3-\nu(A)\wedge e_2.\label{eq:wedge-minus}
\end{equation}
For the standard pair \(M=2^n\), \(A=3^n\), this is
\[
 ne_2\wedge e_3-ne_3\wedge e_2=2n e_2\wedge e_3,
\]
not zero in odd characteristic.  The alternating expression that \emph{does} vanish on the
standard ray is instead
\begin{equation}
 \tau^+_{\wedge}=\nu(M)\wedge e_3+\nu(A)\wedge e_2,\label{eq:wedge-plus}
\end{equation}
which has the wrong sign relative to the real determinant.  This elementary sign reversal is
the core parity barrier.

\subsection{The standard-vanishing \texorpdfstring{\(K_2\)}{K2} class}

By Matsumoto's theorem, Milnor and Quillen \(K_2\) agree for \(\Q\)
\cite{MilnorK,WeibelK}.  Define
\begin{equation}
 \Xi(M,A)=\{M,3\}\,\{A,2\}\in K_2^M(\Q).\label{eq:Xi}
\end{equation}
Bilinearity and antisymmetry give
\[
 \Xi(2^n,3^n)=\{2,3\}^n\{3,2\}^n=1.
\]
Thus \(\Xi\) detects the exterior expression \eqref{eq:wedge-plus}, not the symmetric tensor
\eqref{eq:symtensor}.

For a rational prime \(p\), use the tame-symbol convention
\begin{equation}
 \partial_p\{u,v\}
 =(-1)^{v_p(u)v_p(v)}
 \overline{u^{v_p(v)}v^{-v_p(u)}}
 \in\F_p^\times.\label{eq:tamesymbol}
\end{equation}
Write \(m_p=v_p(M)\) and \(a_p=v_p(A)\).

\begin{proposition}[Exact tame symbols]
\label{prop:tamesymbols}
The class \(\Xi(M,A)\) has the following residues:
\begin{align}
 \partial_p\Xi(M,A)
 &=2^{-a_p}3^{-m_p}\pmod p,
 &&p\notin\{2,3\},\label{eq:tameoutside}\\
 \partial_3\Xi(M,A)
 &=(-1)^{m_3+a_3}\overline{M/3^{m_3}}
 \in\F_3^\times,\label{eq:tame3}
\end{align}
and \(\partial_2\Xi\) is automatically trivial because \(\F_2^\times=\{1\}\).
\end{proposition}

\begin{proof}
For \(p\notin\{2,3\}\), formula \eqref{eq:tamesymbol} gives
\(\partial_p\{M,3\}=3^{-m_p}\) and
\(\partial_p\{A,2\}=2^{-a_p}\).  Multiplication yields
\eqref{eq:tameoutside}.  At \(p=3\),
\[
 \partial_3\{M,3\}=(-1)^{m_3}\overline{M/3^{m_3}},
 \qquad
 \partial_3\{A,2\}=2^{-a_3}=(-1)^{a_3},
\]
which proves \eqref{eq:tame3}.
\end{proof}

\begin{warning}
The symbol \(\Xi\) is not known to be trivial for a counterexample.  The proposition is a
conditional sieve: it records the exact consequences \emph{if} an additional argument can
force \(\Xi=1\).
\end{warning}

\begin{corollary}[Congruence consequences of \(\Xi=1\)]
\label{cor:K2sieve}
If \(\Xi(M,A)=1\), then for every \(p\notin\{2,3\}\),
\begin{equation}
 2^{a_p}3^{m_p}\equiv1\pmod p.\label{eq:K2congruence}
\end{equation}
In particular:
\begin{itemize}
\item if \(p\mid M\) and \(p\nmid A\), then \(\ord_p(3)\mid m_p\);
\item if \(p\mid A\) and \(p\nmid M\), then \(\ord_p(2)\mid a_p\);
\item if \(p\) divides both, then their valuation pair lies on the cyclic congruence
\(2^{a_p}3^{m_p}=1\) in \(\F_p^\times\).
\end{itemize}
\end{corollary}

There is also a converse at the level of all tame symbols.  The localization sequence \cite{WeibelK} gives
\begin{equation}
 0\longrightarrow K_2(\Z)\longrightarrow K_2(\Q)
 \xrightarrow{\oplus_p\partial_p}
 \bigoplus_p\F_p^\times\longrightarrow0,\label{eq:localization}
\end{equation}
with \(K_2(\Z)\cong\Z/2\Z\), generated by \(\{-1,-1\}\).  If every finite tame symbol of
\(\Xi\) is trivial, then \(\Xi\) lies in this order-two kernel.  The real Hilbert-symbol
signature of a symbol built from positive entries is trivial, whereas \(\{-1,-1\}\) has
nontrivial real signature.  Hence \(\Xi=1\).  Thus the congruences above, together with the
prime-three condition, exactly characterize the triviality of this positive \(K_2\) class.

\subsection{Why norm-residue reciprocity does not encode \texorpdfstring{\(Q_{M,A}\)}{Q(M,A)}}

For odd \(q\), the norm-residue isomorphism identifies
\[
 K_2^M(\Q)/q\cong H^2(\Q,\mu_q^{\otimes2}).
\]
The degree-one Kummer cup product is graded-commutative:
\[
 \alpha\smile\beta=-\beta\smile\alpha.
\]
It therefore factors through the exterior square.  In contrast, the finite determinant
\(Q_{M,A}(\chi)\) is obtained by evaluating a \emph{symmetric} quadratic tensor on one
character twice.  Ordinary degree-two Kummer cohomology has no natural operation that
changes the sign in \eqref{eq:wedge-plus} back into the sign of
\(\tau_{\mathrm{sym}}\) while preserving standard-ray vanishing.

\begin{resultbox}[Cohomological parity barrier]
A direct proof by the usual odd-primary \(K_2\) symbols or cup products cannot simply encode
the Alaoglu--Erd\H{o}s determinant: the standard-vanishing exterior class has the opposite
sign, while the same-sign antisymmetrization fails to vanish even for integer exponents.
Any reciprocity proof must introduce genuinely symmetric degree-two information, not only an
ordinary Steinberg symbol.
\end{resultbox}

Characteristic two makes plus and minus coincide, and \Cref{sec:char2} shows that the
resulting Boolean quadratic is visible to Frobenius at positive density.  This still does not
turn the ordinary local Hilbert symbol into \(\Delta_{\ell,2}\): for an odd eligible
\(\ell\), all four radicands are \(\ell\)-adic units, and the Hilbert symbol of two units
over \(\Q_\ell\) is trivial.  Thus the diagonal Boolean Frobenius function is not the local
invariant of \(\Xi\).  Quadratic refinements, Bocksteins, or operations related to the
Pontryagin square remain conceivable research directions, but no construction is known that
connects them to the positive real logarithmic period.

\section{Additional directions tested and their exact stopping points}
\label{sec:attempts}

The new torsion viewpoint suggests several apparently promising shortcuts.  None currently
closes the conjecture; recording their exact failure modes is part of the progress.

\subsection{Sending \texorpdfstring{\(q\)}{q} to infinity}

For generic \(q\), \Cref{thm:genericdensity} makes almost every character in \(W_q^\vee\)
detect the tensor.  This does not force a contradiction because the rational primes that
support \(q\)-th power-residue characters satisfy \(\ell\equiv1\pmod q\), a set of density
\(1/(q-1)\).  One would need a uniform theorem coupling many torsion levels at the same
rational primes.  The fields \(\mathcal L_q\) vary with \(q\), and no independence statement
strong enough for this purpose follows from ordinary Chebotarev alone.

\subsection{Reducing the real exponent modulo \texorpdfstring{\(q\)}{q}}

For an integer exponent \(n\), the identities
\[
 \chi(M)=n\chi(2),\qquad \chi(A)=n\chi(3)
\]
make the determinant vanish.  It is tempting to replace \(n\) by a ``residue class of
\(\beta\).''  There is no canonical map \(\mathbb R\to\F_q\) compatible with
exponentiation on positive rationals.  Choosing rational approximants \(r/s\) merely replaces
this nonexistent residue by the two finite errors
\[
 s\chi(M)-r\chi(2),\qquad s\chi(A)-r\chi(3),
\]
which are unrelated to the smallness of \(s\beta-r\) at the real place.

If one deliberately chooses \(q\mid s\), both errors become scalar multiples of
\(\chi(2),\chi(3)\) and their wedge vanishes for the trivial reason that \(s=0\) in
\(\F_q\).  This produces no information about \(\Delta_{\ell,q}\).  If \(q\nmid s\), one
recovers the common-gap criterion in \Cref{prop:gapprime}, but only at primes dividing both
integer gaps.

\subsection{Global reciprocity without an auxiliary object}

Artin reciprocity classifies abelian Frobenius data, and the norm-residue theorem identifies
Steinberg symbols with cup products.  Neither theorem supplies an algebraic object whose
archimedean logarithm is
\(\log M\log3-\log A\log2\).  Reciprocity can compare local invariants of an already
constructed class; it does not manufacture the missing symmetric class.  The sign computation
in \Cref{sec:K2} shows that the most obvious candidate lives in the wrong functor.

\subsection{Forcing a small detecting prime unconditionally}

\Cref{thm:effectivewitness} gives polynomial bounds in \(MA\) without GRH and
polylogarithmic bounds under GRH at both the quadratic and cubic levels.  Even a uniform
numerical bound such as \(\ell<10^6\) would
not prove the conjecture: a detecting prime is expected for a counterexample.  The needed
step is not to find a nonzero determinant, but to contradict it using the real equation.

\subsection{Treating the density theorem as probabilistic evidence}

The Chebotarev density is exact, not a random-model prediction.  Conversely, numerical
agreement with that density cannot distinguish a near miss from an exact counterexample.
The field \(\mathcal L_q\) contains no archimedean memory of how small
\(Q((\log p)_p)\) is.  The experiment in \Cref{sec:computation} is therefore best read as a
no-go result for continuity heuristics.

\begin{center}
\begin{tabularx}{0.96\textwidth}{@{}p{0.25\textwidth}X@{}}
\toprule
Attempt & Exact obstruction \\
\midrule
Large odd torsion level \(q\) & Relative density rises, but eligible rational primes have density
\(1/(q-1)\). \\
A residue class for \(\beta\) & No topology-compatible homomorphism
\(\mathbb R\to\F_q\) exists. \\
Continued fractions & Real smallness controls gap size, not common Frobenius classes. \\
Ordinary \(K_2\) reciprocity & Alternating cup products have the wrong sign symmetry. \\
Effective witness search & A witness confirms nonstandardness; it does not contradict the
real relation. \\
Finite experiments & Exact real equality is an infinite transcendence condition. \\
\bottomrule
\end{tabularx}
\end{center}

\section{A focused research program}
\label{sec:program}

The new results narrow the next useful work to four programs.  They are ordered by how
directly they attack \Cref{crit:threshold}.

\subsection{Program I: Chebotarev-weighted common-divisor theory}

For convergents \(r_n/s_n\to\beta\), study
\[
 G_n=\gcd(M^{s_n}-2^{r_n},A^{s_n}-3^{r_n})
\]
not merely through \(\log G_n\), but through its Frobenius support.  For odd \(s_n\),
\Cref{cor:forbidden} excludes the quadratic detecting classes; for \(3\nmid s_n\), it
excludes the cubic detecting classes.  A theorem showing that synchronized near-equalities
force prime divisors to sample all Kummer conjugacy classes with a positive lower frequency
would contradict these restrictions.

A concrete target is the following.

\begin{gapbox}[Prime-sampling target]
Show that for infinitely many convergents with \(3\nmid s_n\), the common divisor \(G_n\)
has at least one prime factor in every fixed positive-density Chebotarev set attached to the
four radicands \(2,3,M,A\), provided \(G_n\) exceeds an explicit height threshold.
\end{gapbox}

This statement is far stronger than standard results on primitive divisors of one recurrence;
it concerns a gcd of two multiplicatively independent exponential sequences with synchronized
irrational slopes.  The formulation makes the required strengthening explicit.

\subsection{Program II: a symmetric reciprocity object}

The desired finite function is
\[
 \chi\longmapsto Q_{M,A}(\chi),
\]
a quadratic function on \(H^1\), not an alternating bilinear cup product.  Possible homes for
such information include divided-power algebras, quadratic refinements at the prime two,
secondary cohomology operations, or regulators on a higher motivic object.  Any candidate
must pass three tests:
\begin{enumerate}
\item it vanishes identically for \((M,A)=(2^n,3^n)\);
\item its finite localizations recover \(\Delta_{\ell,q}\) or a function with the same zero
locus; and
\item its real regulator contains \(\log M\log3-\log A\log2\), not an alternating wedge.
\end{enumerate}
The ordinary Steinberg symbol fails the third test by \Cref{sec:K2}.

\subsection{Program III: a character-sensitive auxiliary determinant}

The existing report's auxiliary-function constructions are insensitive to Kummer conjugacy
classes.  A new determinant could incorporate the idempotent of a detecting orbit in the
group algebra
\[
 \Q(\zeta_q)[\Gal(\mathcal L_q/K_q)].
\]
The desired architecture is:
\begin{enumerate}
\item choose a character \(\chi\) with \(Q(\chi)\ne0\);
\item project an algebraic auxiliary expression to the corresponding isotypic component;
\item use the exact real period zero to make that component unusually small at one real
embedding; and
\item use integrality and product-formula control at the other embeddings to force it to
vanish, contradicting the character choice.
\end{enumerate}
The obstacle is simultaneous control of all conjugates.  Crude norms multiply too many
uncontrolled terms, exactly as in standard radical auxiliary constructions.  The group-ring
projection is the new ingredient worth testing.

\subsection{Program IV: exact formal and computational infrastructure}

The finite algebra in this report is particularly suitable for formalization and exhaustive
checking.  A reusable implementation should:
\begin{itemize}
\item represent prime exponent vectors as finitely supported integer maps;
\item compute \(Q_{M,A}\), its content, and its rank over \(\Q\) and \(\F_q\);
\item verify a power-residue certificate \((q,\ell,\omega_\ell)\);
\item enumerate the finite character space and check the exact density numerator; and
\item export a small certificate independent of integer factorization heuristics.
\end{itemize}
This will not prove the transcendence bridge, but it removes all finite ambiguity from future
candidate arguments.

\section{Lean formalization blueprint}
\label{sec:lean}

This section describes a realistic formalization boundary.  It does not claim that the new
theorems have already been checked by Lean.

\subsection{Finite algebra layer}

The following components should be routine in Mathlib once the attachment's primitive-output
objects are imported.

\begin{enumerate}
\item \textbf{Prime exponent vectors.}  Use a finitely supported function
\(\mathbb P\to_0\Z\) for \(\nu(u)\), together with its reduction modulo \(q\).
\item \textbf{Quadratic obstruction.}  Rather than implement an infinite polynomial ring,
represent
\[
 Q(z)=\langle\nu(M),z\rangle z(3)-\langle\nu(A),z\rangle z(2)
\]
as a function on finitely supported vectors.  A finite support containing
\(\Supp(6MA)\) suffices.
\item \textbf{Content theorem.}  Expand using \eqref{eq:Qexpanded}; prove that the gcd of the
three coefficient blocks is \(\gcd(d,e,a-b)\).  The central interface with \ReportA{} is
the theorem
\[
 \texttt{simultaneousCoordinateGCD\_eq\_one}.
\]
\item \textbf{Nonzero reduction.}  A primitive finite coefficient family has a nonzero member
modulo every prime.
\item \textbf{Quadratic zero count.}  Either formalize the rank classification over finite
fields of odd characteristic, or prove the weaker bound \eqref{eq:quadzerobound} by a direct
fiber argument and reserve the exact rank formulas for a later file.
\item \textbf{Certificate checker.}  Verify that \(\omega^q=1\), \(\omega\ne1\), compute the
four character values by modular exponentiation, and check the determinant.
\end{enumerate}

Schematic theorem signatures, deliberately avoiding unverified Mathlib API names, are:
\begin{verbatim}
content_twoStar (normalization : SimultaneousPrimitiveOutput M A) :
  content (twoStarCoefficients M A) = gcd normalization.d
    (gcd normalization.e |normalization.a - normalization.b|)

primitive_twoStar_of_least (h : IsLeastNonintegerSolution beta) :
  Primitive (twoStarCoefficients (2^beta) (3^beta))

quadratic_nonzero_count (q : Nat.Prime) (hq : q != 2)
    (Q : QuadraticForm (ZMod q) W) (hQ : Q != 0) :
  ((Fintype.card W - #{x | Q x = 0}) : Rat) / Fintype.card W
    >= (1 - 1/q)^2
\end{verbatim}
These are specifications, not compilable declarations.

\subsection{Classical number-theory layer}

Kummer theory and Chebotarev are not presently small local lemmas.  A practical formal report
can separate them as named classical inputs:
\begin{itemize}
\item Kummer duality for a finite-dimensional subgroup of \(K^\times/(K^\times)^q\);
\item the Frobenius/power-residue-symbol correspondence;
\item Chebotarev density for a conjugacy-stable finite subset; and
\item effective Chebotarev under GRH.
\end{itemize}
A first formal milestone can prove the finite statement
\[
 Q\ne0\quad\Longrightarrow\quad
 \#\{\chi:Q(\chi)\ne0\}\ge(q-1)^2q^{r-2},
\]
and then state the density theorem conditionally on an abstract equidistribution axiom.  This
would keep the machine-checked and classical layers visibly distinct.

\subsection{Certificate format}

For \(q=3\), a compact witness consists of
\[
 (M,A,\ell,\omega; c_2,c_3,c_M,c_A),
\]
with checks
\begin{align*}
 &\ell\text{ prime},\quad \ell\equiv1\pmod3,
 \quad \ell\nmid6MA,\\
 &\omega^3\equiv1\pmod\ell,\quad \omega\not\equiv1\pmod\ell,\\
 &u^{(\ell-1)/3}\equiv\omega^{c_u}\pmod\ell
 \quad(u=2,3,M,A),\\
 &c_Mc_3-c_Ac_2\not\equiv0\pmod3.
\end{align*}
No discrete-log search is needed in the checker: the exponents \(c_u\in\{0,1,2\}\) are part
of the certificate.

\section{Audit of the claimed August 2026 breakthroughs}
\label{sec:news}

The surrounding claims were checked only to establish context; they are not mathematical
inputs to any theorem above.

\begin{enumerate}
\item \textbf{Jacobian conjecture.}  A public 2026 preprint by Shuhong Gao gives a
self-contained account of counterexamples in every dimension greater than two, attributing
the initial three-dimensional counterexample to Levent Alp\"oge and the AI-assisted discovery.
An independent Lean 4 verification of the explicit three-dimensional map is publicly archived.
The plane Jacobian conjecture remains outside that result \cite{Gao2026,LeanJacobian}.
\item \textbf{Elliptic curve rank at least 30.}  ICARM's Elliptic Curve Rank Leaderboard lists
curve \#273 with 30 explicit independent rational points and hence a certified lower bound
\(\rank E(\Q)\ge30\), submitted on 20 August 2026.  The page's commentary attributes the
search to Claude with Levent Alp\"oge and Ava Howell.  ``Exactly 30'' is stated there only
under GRH and BSD; the unconditional public certificate is the lower bound \cite{ICARM30}.
\item \textbf{Hadamard orders below 2000.}  Epoch AI reports, provisionally, that a puzzle
posted by a team of three humans and Claude encodes matrices for the twelve previously unknown
admissible orders below 2000, including 668.  OEIS records the same twelve orders as found in
2026.  The attribution and degree of independent reproduction remain more provisional than
the existence claim \cite{EpochHadamard,OEISHadamard}.
\item \textbf{Alaoglu--Erd\H{o}s.}  Targeted searches on 22 August 2026 found no public
preprint, formal certificate, or detailed announcement of the private breakthrough described
in the prompt.  Absence from a search index is not evidence that no private result exists.
\end{enumerate}

The relevant public pages are recorded in the bibliography.  The present report was developed
independently from the attached mathematical corpus and those public sources; no alleged
private argument was available to inspect.

\section{Conclusion}
\label{sec:conclusion}

The real equation of a hypothetical counterexample annihilates a primitive symmetric
quadratic tensor.  The main new observation is that the attachment's affine-primitivity gcd is
exactly the coefficient content of that tensor.  Once this is recognized, Kummer theory makes
the tensor visible at every odd torsion level, and Chebotarev forces nonvanishing at a uniform
positive density of rational primes.

The two smallest torsion levels are especially clean:
\[
 \boxed{\text{least counterexample}\quad\Longrightarrow\quad
 \dens\{\ell:\Delta_{\ell,2}\ne0\}\ge\frac14,
 \qquad
 \dens\{\ell:\Delta_{\ell,3}\ne0\}\ge\frac29.}
\]
The fields producing these witnesses have degrees at most 16 and 162.  Under GRH, their first
witnesses are polylogarithmic in \(MA\).  These statements are exact and effectively
testable.

The same theorem also clarifies why it is not yet a proof.  Finite torsion is generically
transverse; a counterexample would not look locally singular.  What is missing is a theorem
that makes the exact real rank-one relation force too many finite rank drops.  The robust cubic threshold is now explicit: more than \(5/9\) zero density among
\(\ell\equiv1\pmod3\) would suffice.  The quadratic level supplies the stronger absolute
certificate density, but its complementary zero threshold \(3/4\) is less demanding only as
a certificate and more demanding as a transfer theorem.

The algebraic \(K\)-theory audit reaches a similarly sharp conclusion.  Tame symbols give a
valid local congruence sieve, but ordinary cup products are alternating, while the period
obstruction is symmetric.  The obvious reciprocity class either has the wrong sign or fails
to vanish on the standard integer ray.  A successful cohomological proof therefore needs a
symmetric quadratic refinement, not merely a Steinberg symbol.

\begin{gapbox}[Most promising next step]
Combine the positive-density detecting set with the synchronized gap sequence
\(\gcd(M^s-2^r,A^s-3^r)\).  The new sieve says every common prime divisor lies on a finite
Kummer quadric whenever the chosen torsion prime \(q\) does not divide \(s\).  A prime-distribution theorem forcing these gcds to
sample the complementary Chebotarev classes would convert the present transversality theorem
into a contradiction.
\end{gapbox}

\appendix

\section{Reproducibility code for the cubic experiment}
\label{app:code}

The following compact Python program reproduces the character calculation and empirical
density check for a supplied list of pairs.  It requires \texttt{sympy}.  The exhaustive
near-miss search used the same routines and checked both nearest integers to \(M^\theta\).

\begin{lstlisting}[style=pythonstyle]
from itertools import product
from sympy import factorint, primerange, primitive_root

Q = 3
PAIRS = [
    (9329, 1959024), (8427, 1667418), (10855, 2490703),
    (9566, 2038489), (10506, 2364980), (9517, 2021964),
]

def q_character(u: int, ell: int, q: int, primitive: int) -> int:
    omega = pow(primitive, (ell - 1) // q, ell)
    value = pow(u % ell, (ell - 1) // q, ell)
    cur = 1
    for exponent in range(q):
        if cur == value:
            return exponent
        cur = (cur * omega) % ell
    raise RuntimeError("power-residue value not in mu_q")

def q_polynomial(M: int, A: int, support, z, q: int) -> int:
    fm, fa = factorint(M), factorint(A)
    LM = sum(fm.get(p, 0) * z[i] for i, p in enumerate(support)) % q
    LA = sum(fa.get(p, 0) * z[i] for i, p in enumerate(support)) % q
    X2 = z[support.index(2)]
    X3 = z[support.index(3)]
    return (LM * X3 - LA * X2) % q

for M, A in PAIRS:
    support = sorted(set(factorint(6 * M * A)))

    # Exact finite-character proportion.  Restriction to W_q has equal fibers.
    total = q ** len(support)
    nonzero = sum(
        q_polynomial(M, A, support, z, q) != 0
        for z in product(range(q), repeat=len(support))
    )

    eligible = detected = 0
    first = None
    for ell in primerange(2, 200001):
        if ell in support or (ell - 1) % q != 0:
            continue
        primitive = primitive_root(ell)
        c2 = q_character(2, ell, q, primitive)
        c3 = q_character(3, ell, q, primitive)
        cM = q_character(M, ell, q, primitive)
        cA = q_character(A, ell, q, primitive)
        delta = (cM * c3 - cA * c2) % q
        eligible += 1
        if delta:
            detected += 1
            if first is None:
                first = (ell, delta, (c2, c3, cM, cA))

    print(
        M, A,
        f"exact={nonzero}/{total}={nonzero/total:.6f}",
        f"empirical={detected}/{eligible}={detected/eligible:.6f}",
        f"first={first}",
    )
\end{lstlisting}

For the six pairs in \Cref{comp:table}, the exact finite-character proportion is \(16/27\).
The prime cutoff \(200000\) gives approximately 8990 eligible primes per pair.

\section{Proof details for the finite-field counts}
\label{app:counts}

For completeness, let \(Q\) have rank \(\rho\) on an \(r\)-dimensional vector space over
\(\F_q\), \(q\) odd.  Factoring out the radical contributes a multiplicative factor
\(q^{r-\rho}\) to the number of zeros.  For a nondegenerate form in \(\rho\) variables,
\[
 N_0(Q)=
 \begin{cases}
 q^{\rho-1},&\rho\text{ odd},\\[0.3em]
 q^{\rho-1}+\varepsilon(q-1)q^{\rho/2-1},&\rho\text{ even},
 \end{cases}
\]
where \(\varepsilon=1\) for the split even-dimensional form and \(-1\) for the nonsplit
form.  At rank two, the split count is \(2q-1\), the maximum among all nonzero quadratic
forms.  Multiplying by the radical factor proves \eqref{eq:quadzerobound}.

For the rank-four hyperbolic form \(z_1z_2-z_3z_4\), the formula gives
\[
 N_0=q^3+(q-1)q,
 \qquad
 1-\frac{N_0}{q^4}=1-\frac1q-\frac{q-1}{q^3}.
\]
For \(q=3\), this is \(16/27\), the exact value used in the computational table.

\section{Source and dependency ledger}
\label{app:ledger}

\begin{longtable}{@{}>{\raggedright\arraybackslash}p{0.29\textwidth}>{\raggedright\arraybackslash}p{0.20\textwidth}>{\raggedright\arraybackslash}p{0.42\textwidth}@{}}
\toprule
Statement & Status & Dependencies \\
\midrule
\endfirsthead
\toprule
Statement & Status & Dependencies \\
\midrule
\endhead
Least-output normalization and gcd one & inherited & Attached \ReportA, simultaneous primitive output normalization \\
Rank three/four of the symmetric tensor & inherited/classical & Attached \ReportA; Baker linear independence \\
Content formula \eqref{eq:content} & proved here & Elementary coefficient comparison \\
Kummer rank injection & proved here & Norm argument; standard Kummer theorem \\
Exact density \eqref{eq:absdensity} & proved here from classical input & Kummer--Frobenius correspondence; Chebotarev \\
Universal odd-primary density \eqref{eq:absolutebound} & proved here & Content formula; finite quadratic zero count; Chebotarev \\
Quadratic density \eqref{eq:quadraticdensity} & proved here & Content formula; Boolean minimum distance; multiquadratic Chebotarev \\
Generic exact densities & proved here & Inherited rank classification; finite-field counts \\
GRH low-torsion witness bounds & conditional & Effective Chebotarev under GRH; fixed-degree discriminant bounds \\
Unconditional witness bounds & proved from classical effective theorem & Unconditional least-prime Chebotarev bound \\
Gap-prime sieve & proved here & Elementary character calculation \\
Tame-symbol formulas & proved here & Matsumoto theorem; localization sequence for \(K_2(\Q)\) \\
Parity barrier & proved here as a structural no-go & Graded commutativity of cup products \\
Near-miss table & computation & Python/Sympy exact modular arithmetic \\
Cubic density transfer & open target & Any archimedean-to-Kummer rank-drop theorem \\
\bottomrule
\end{longtable}

\begin{thebibliography}{99}

\bibitem{AE1944}
L. Alaoglu and P. Erd\H{o}s,
\emph{On highly composite and similar numbers},
Transactions of the American Mathematical Society \textbf{56} (1944), 448--469.
\url{https://doi.org/10.1090/S0002-9947-1944-0011087-2}.

\bibitem{ReportA}
\emph{The Alaoglu--Erd\H{o}s Integer-Exponent Conjecture: Machine-Checked Partial Results,
Exact Conditional Structure, Determinant Methods, and the Remaining Barrier},
attached unified report A, 20,925 lines, accessed 22 August 2026.
Decompressed SHA-256:
\texttt{6bb33ec36c1d548ae8f343681bda054e752d62a24a0ac4157e4471dbd0a5c565}.

\bibitem{Baker1966}
A. Baker,
\emph{Linear forms in the logarithms of algebraic numbers},
Mathematika \textbf{13} (1966), 204--216.

\bibitem{Lang1966}
S. Lang,
\emph{Introduction to Transcendental Numbers},
Addison--Wesley, 1966.

\bibitem{Ramachandra}
K. Ramachandra,
\emph{Contributions to the theory of transcendental numbers, I--II},
Acta Arithmetica \textbf{14} (1967/68), 65--88 and 73--88.

\bibitem{MilneFT}
J. S. Milne,
\emph{Fields and Galois Theory}, version 5.10, 2022,
\url{https://www.jmilne.org/math/CourseNotes/ft.html}.

\bibitem{MilneANT}
J. S. Milne,
\emph{Algebraic Number Theory}, version 3.08, 2020,
\url{https://www.jmilne.org/math/CourseNotes/ant.html}.

\bibitem{MilneCFT}
J. S. Milne,
\emph{Class Field Theory}, version 4.03, 2020,
\url{https://www.jmilne.org/math/CourseNotes/cft.html}.

\bibitem{Neukirch}
J. Neukirch,
\emph{Algebraic Number Theory},
Springer, 1999.

\bibitem{SerreCheb}
J.-P. Serre,
\emph{Lectures on \(N_X(p)\)},
Research Notes in Mathematics 11, CRC Press, 2012.

\bibitem{LagariasOdlyzko}
J. C. Lagarias and A. M. Odlyzko,
\emph{Effective versions of the Chebotarev density theorem},
in \emph{Algebraic Number Fields}, Academic Press, 1977, 409--464.

\bibitem{GrenieMolteni}
L. Greni\'e and G. Molteni,
\emph{An effective Chebotarev density theorem under GRH},
Journal of Number Theory \textbf{200} (2019), 441--485;
\url{https://arxiv.org/abs/1709.07609}.

\bibitem{ThornerZaman}
J. Thorner and A. Zaman,
\emph{A unified and improved Chebotarev density theorem},
Algebra \& Number Theory \textbf{13} (2019), 1039--1068;
\url{https://arxiv.org/abs/1803.02823}.

\bibitem{KadiriWong}
H. Kadiri and P.-J. Wong,
\emph{Primes in the Chebotarev density theorem for all number fields},
Journal of Number Theory \textbf{241} (2022), 550--582;
\url{https://arxiv.org/abs/2105.14181}.

\bibitem{MilnorK}
J. Milnor,
\emph{Introduction to Algebraic \(K\)-Theory},
Annals of Mathematics Studies 72, Princeton University Press, 1971.

\bibitem{WeibelK}
C. A. Weibel,
\emph{The \(K\)-Book: An Introduction to Algebraic \(K\)-Theory},
Graduate Studies in Mathematics 145, American Mathematical Society, 2013.

\bibitem{Gao2026}
S. Gao,
\emph{Counterexamples to the Jacobian conjecture in dimensions greater than two},
arXiv:2608.00222, 2026;
\url{https://arxiv.org/abs/2608.00222}.

\bibitem{LeanJacobian}
P. Nogueira Grossi,
\emph{An Independent Lean 4 Verification of the Alp\"oge--Fable Counterexample},
Zenodo, 23 July 2026;
\url{https://doi.org/10.5281/zenodo.21514514}.

\bibitem{ICARM30}
Institute for Computer-Aided Reasoning in Mathematics,
\emph{Elliptic Curve Rank Leaderboard, curve \#273},
accessed 22 August 2026;
\url{https://elliptic-rank.icarm.cloud/curve/273}.

\bibitem{EpochHadamard}
Epoch AI,
\emph{Hadamard Matrix of Order 668},
provisional solution update, accessed 22 August 2026;
\url{https://epoch.ai/frontiermath/open-problems/hadamard}.

\bibitem{OEISHadamard}
OEIS Foundation,
\emph{A019442: Numbers \(m\) such that a Hadamard matrix of order \(m\) exists},
2026 update;
\url{https://oeis.org/A019442}.

\end{thebibliography}

\end{document}
